\documentclass[11pt,a4paper]{article}
\usepackage[UTF8]{ctex}
\usepackage{amsmath, amssymb, amsthm, amsfonts}
\usepackage{geometry}
\usepackage{enumitem}

\geometry{left=2cm, right=2cm, top=2.5cm, bottom=2.5cm}

\begin{document}

\section*{集合与不等式 答案解析}

\begin{enumerate}[label=\textbf{题 \arabic*}]

\item \textbf{答案：} C

\textbf{解析：} 由 \( x^2 - 3x - 4 < 0 \) 解得 \( (x-4)(x+1) < 0 \)，即 \( -1 < x < 4 \)，
所以 \( B = (-1, 4) \)。又 \( A = [-2, 5] \)，因此 \( A \cup B = [-2, 5] \)，故选 C。

\item \textbf{答案：} C

\textbf{解析：} 因为指数函数 \( y = 2^x \) 在 \( \mathbb{R} \) 上单调递增，由 \( a > b \) 可得 \( 2^a > 2^b \)，故 C 正确。
A 反例：\( a = -1, b = -2 \) 时 \( a^2 = 1 < 4 = b^2 \)；
B 反例：\( a = 2, b = -1 \) 时 \( \frac{1}{2} > -1 \)；
D 要求 \( a, b > 0 \) 才成立。

\item \textbf{答案：} \( [-2, 1) \cup (1, +\infty) \)

\textbf{解析：} 由 \( x+2 \geq 0 \) 得 \( x \geq -2 \)，由分母 \( x-1 \neq 0 \) 得 \( x \neq 1 \)，
因此定义域为 \( [-2, 1) \cup (1, +\infty) \)。

\end{enumerate}

\end{document}
